Let \(c_{\epsilon}^{\mathrm L}>0\) be the constant supplied by \(\mathrm{thm:all\mbox{-}estimator\mbox{-}lower}\), and let \(C_{\epsilon}^{\mathrm U}<\infty\) be the constant supplied by \(\mathrm{thm:jackson\mbox{-}factorial\mbox{-}upper}\).  These two established results apply for every \(n\geq1\) and \(d\geq2\) under the common observed-law class \(\mathcal D_{d,\epsilon}^{\mathrm{obs}}\).  Hence
\[
 \mathfrak R_{n,d,\epsilon}
 \geq c_{\epsilon}^{\mathrm L}
 \min\!\left\{1,\frac{d}{n\log(ed)}\right\},                                      \tag{1}
\]
and
\[
 \sup_{\mathbb P\in\mathcal D_{d,\epsilon}^{\mathrm{obs}}}
 E_{\mathbb P}\!\left[
 \left\{\widehat V_{n,d}^{\mathrm{JF}}-\Psi(\mathbb P)\right\}^{2}
 \right]
 \leq C_{\epsilon}^{\mathrm U}
 \min\!\left\{1,\frac{d}{n\log(ed)}\right\}.                                      \tag{2}
\]

It remains to justify the middle inequality rather than identify it with either bound.  By \(\mathrm{def:minimax\mbox{-}risk}\), \(\mathfrak R_{n,d,\epsilon}\) is the infimum, over all measurable sample maps, of the displayed worst-case squared risk.  The explicit all-data estimator whose admissibility and risk are asserted by \(\mathrm{thm:jackson\mbox{-}factorial\mbox{-}upper}\) is one such map.  Evaluating an infimum at this particular admissible estimator therefore gives
\[
 \mathfrak R_{n,d,\epsilon}
 \leq
 \sup_{\mathbb P\in\mathcal D_{d,\epsilon}^{\mathrm{obs}}}
 E_{\mathbb P}\!\left[
 \left\{\widehat V_{n,d}^{\mathrm{JF}}-\Psi(\mathbb P)\right\}^{2}
 \right].                                                                            \tag{3}
\]
Set \(c_\epsilon=c_{\epsilon}^{\mathrm L}\) and \(C_\epsilon=\max\{C_{\epsilon}^{\mathrm U},c_{\epsilon}^{\mathrm L}\}\).  Then \(0<c_\epsilon\leq C_\epsilon<\infty\), and concatenating \((1)\), \((3)\), and \((2)\) proves the three inequalities in the theorem.

Finally, the causal assertion is a transfer, not a definition of the causal target.  Proposition \(\mathrm{prop:identification\mbox{-}and\mbox{-}extension}\) proves both inclusions needed for the experiment: every \(P\in\mathcal P_{d,\epsilon}\) has \(\operatorname{Obs}(P)\in\mathcal D_{d,\epsilon}^{\mathrm{obs}}\), and every \(\mathbb P\in\mathcal D_{d,\epsilon}^{\mathrm{obs}}\) has at least one completion \(P\in\mathcal P_{d,\epsilon}\) with \(\operatorname{Obs}(P)=\mathbb P\).  Proposition \(\mathrm{prop:causal\mbox{-}optimal\mbox{-}value\mbox{-}corollary}\), using consistency and armwise conditional exchangeability, proves for every such completion that
\[
 V^\star(P)=\Psi\{\operatorname{Obs}(P)\}.
\]
Therefore, for every observed-data estimator \(\widehat V\),
\[
 \sup_{P\in\mathcal P_{d,\epsilon}}
 E_{\operatorname{Obs}(P)}\!\left[
 \left\{\widehat V-V^\star(P)\right\}^{2}
 \right]
 =
 \sup_{\mathbb P\in\mathcal D_{d,\epsilon}^{\mathrm{obs}}}
 E_{\mathbb P}\!\left[
 \left\{\widehat V-\Psi(\mathbb P)\right\}^{2}
 \right].                                                                             \tag{4}
\]
Indeed, the first inclusion gives the inequality from left to right in \((4)\), while the completion of each \(\mathbb P\) and the target equality give the reverse inequality.  Taking the infimum over the same measurable observed-data maps on both sides of \((4)\) transfers the already proved statistical frontier to \(V^\star(P)\), which establishes the final sentence.